\SourceLeaf{332}

\[
(4), (8), (9) \text{ from I}
\]
\[
(3), (7), (11) \text{ from II}
\]
\[
(5), (6), (10) \text{ from I and II combined.}
\]

Any of equations (3)—(11) may be used to solve the \textsc{Pothenot} problem, since the two unknowns occurring in each equation have a known sum.

From Figure II the quadrilateral equations

\[
\frac{\sin(\mathfrak{A}-A)}{a\alpha} = \frac{\sin(\mathfrak{B}-B)}{b\beta} = \frac{\sin(\mathfrak{C}-C)}{c\gamma} = \frac{rr-\rho\rho}{2\alpha\beta\gamma r}
\]

follow immediately if $r$ is the radius of the circle through $A$, $B$, $C$ and $\rho$ the distance of the fourth point from the centre of this circle. Equality of the signs of

\[
\sin(\mathfrak{A}-A), \qquad \sin(\mathfrak{B}-B), \qquad \sin(\mathfrak{C}-C)
\]

is the criterion for the physical possibility of the problem.

The geometrical meaning of formulas (2)—(11) can be associated very neatly with the five circles described respectively\ through three points each of Fig.\ I or of Fig.\ II, and whose diameters may be $p$, $q$, $r$, $s$, $t$. We then have

\[
\sin A = \frac{a}{p}, \qquad \sin B = \frac{b}{p}, \qquad \sin C = \frac{c}{p},
\]
\[
\sin\mathfrak{A} = \frac{a}{q}, \qquad \sin 4 = \frac{\gamma}{q}, \qquad \sin 5 = \frac{\beta}{q},
\]
\[
\sin\mathfrak{B} = \frac{b}{r}, \qquad \sin 6 = \frac{\alpha}{r}, \qquad \sin 1 = \frac{\gamma}{r},
\]
\[
\sin\mathfrak{C} = \frac{c}{s}, \qquad \sin 2 = \frac{\beta}{s}, \qquad \sin 3 = \frac{\alpha}{s},
\]
\[
\sin(\mathfrak{A}-A) = \frac{a\alpha}{t}, \qquad \sin(\mathfrak{B}-B) = \frac{b\beta}{t}, \qquad \sin(\mathfrak{C}-C) = \frac{c\gamma}{t}.
\]

Substituting these expressions for the 15 sines in the 10 equations (2)—(11) makes them identities.

\Needspace{7\baselineskip}Just as from equation (2) the others (3)—(11) follow identically, so all the other 9 can be derived from any one of them. E.\,g., one need only put equation (3) in the form
